% Plan de tesis -- Capitulo 3.
%
% Correlato en el plan de ejemplo: capitulo 3 (Experimentos preliminares y resultados), ~2 paginas.
% El ejemplo muestra una tabla chica con la comparacion contra el metodo base y enumera los
% experimentos siguientes. Es el capitulo que convierte la propuesta en algo ya empezado.
%
% Material disponible: la campana d303-a6262712 del repositorio (rama feature/fase-estrategias)
% tiene las celdas necesarias. Los numeros se citan por campana/celda/commit, nunca por
% "Decision NNN" (la numeracion 039-051 esta duplicada entre ramas; ver manuscritos/PLAN.md).
\chapter{Experimentos preliminares y resultados}

% PENDIENTE: redactar a partir de la campana d303-a6262712.
% Lo que ya esta verificado contra el ledger y los JSON de la campana:
%   - ancla B0.06: 0,4948 +/- 0,0098, con una mejora del 7,6 % sobre la linea de base;
%   - celda B1.06, tau fijo: empata con el ancla, no la supera;
%   - contraste de Diebold-Mariano contra persistencia: p = 0,16, es decir que con el periodo de
%     evaluacion disponible la diferencia todavia no es distinguible del azar.
% Ese ultimo punto es el mas importante del capitulo y conviene no maquillarlo: motiva el trabajo
% en lugar de debilitarlo, porque muestra que el problema no esta resuelto.

% PENDIENTE: tabla de resultados preliminares, al estilo de la Tab. 3.1 del plan de ejemplo.

% PENDIENTE: enumerar los experimentos siguientes (modulacion de tau, comparacion contra las
% funciones objetivo que devuelva el relevamiento, evaluacion por regimen).
